%% qiml26.tex — AAAI QIML 2026 short paper
%% Build: pdflatex qiml26.tex && bibtex qiml26 && pdflatex qiml26.tex && pdflatex qiml26.tex
%% Deadline: 2026-09-01 | 6-8 pages excluding references | two-column

\documentclass[letterpaper]{article}
\usepackage[]{aaai2026}
\usepackage{times}
\usepackage{helvet}
\usepackage{courier}
\usepackage[hyphens]{url}
\usepackage{graphicx}
\urlstyle{rm}
\def\UrlFont{\rm}
\usepackage{natbib}
\usepackage{caption}
\usepackage{amsmath,amssymb}
\usepackage{booktabs}
\frenchspacing
\setlength{\pdfpagewidth}{8.5in}
\setlength{\pdfpageheight}{11in}

\pdfinfo{
/TemplateVersion (2026.1)
/Title (The Pauli Lightcone Frontier as a Space-Time Observable for Quantum Noise Characterization)
/Author (Paolo D'Alberto)
/Subject (Quantum Computing, Tensor Networks, Noise Characterization)
/Keywords (Pauli lightcone, operator spreading, noise characterization, multi-product formula, tensor network)
}

% ─────────────────────────────────────────────────────────────────────────────

\title{The Pauli Lightcone Frontier as a Space-Time Observable\\
       for Quantum Noise Characterization}

\author{Paolo D'Alberto}
\affiliations{
Advanced Micro Devices, Inc.\\
paolo.dalberto@amd.com%
}

% ─────────────────────────────────────────────────────────────────────────────
\begin{document}
\maketitle

\begin{abstract}
We simulate the kicked Heisenberg Ising model on two lattices and
study the non-identity Pauli weight $n(v,t)$, the probability that the
evolved operator has support at site $v$ at cycle $t$, as a space-time
observable.  Under noise amplification at levels $\gamma$, the
$n(v,t)$ directly reveals how noise deforms the operator lightcone
across the lattices.  We introduce the \emph{wavemap}: a spatial
portrait of noise effects assigning each site a per-noise-level
arrival delay $l_\gamma(v)$ and cross-entropy loss
$\mathcal{L}_\gamma(v)$.  The wavemap provides direct observable
signatures of the noise eigenvalue structure: pure amplitude damping
produces a uniform spatial delay with no weight injection outside the
lightcone, while rotation-type noise injects weight at unvisited
sites, making the frontier appear locally faster.  Applying
the multi-product formula (MPF) with time-adaptive coefficients
$\alpha(t)$ over the frontier recovers up to 55\% of the information
loss relative to the best noisy sample.  The noise amplification
parameter $\gamma$ maps directly to the gate-folding factor in
zero-noise extrapolation, making the wavemap methodology immediately
applicable to real quantum processors.
\end{abstract}

% ─────────────────────────────────────────────────────────────────────────────
\section{Introduction and Setup}
\label{sec:intro}

We simulate the kicked isotropic Heisenberg model~\cite{Begusic2025RealTime}:
\begin{equation}
H = J\!\sum_{\langle ij\rangle}(X_iX_j+Y_iY_j+Z_iZ_j),\quad
U_{\rm kick}=\textstyle\prod_v R_x(2\phi)
\end{equation}
with $\phi=1.414$,
on two topologies: a heavy-hex $3\times3$ lattice (68 sites,
76 bonds) with per-bond noise PTM representative of realistic
heavy-hex device noise (uniform $J=1$ couplings),
and a rectangular $7\times7$ lattice (49 sites, 84 bonds) with synthetic
uniform depolarizing noise.
We use a GPU-accelerated belief propagation tensor network simulator~\cite{dalberto2026ising}
with Trotter step $\varepsilon=0.10$, bond dimension $\chi$, and
truncation cutoff $\delta=10^{-3}$ ($\chi=200$ heavy-hex, $\chi=50$
rect $7\times7$). The simulator is a GPU-first implementation of the
belief propagation tensor network methodology; the Julia reference
implementation used in~\cite{robertson2026lindbladmultiproductformulas}
is based on ITensor~\cite{ITensor} and
TensorNetworkQuantumSimulator~\cite{tindall2023gauging_bp,rudolph2025simulating2d}.

Noise enters via Pauli transfer matrices
(PTM)~\cite{rudolph2025simulating}: each bond PTM composes the ideal
gate with $M_\gamma^{(b)} = \gamma M_{\rm hw}^{(b)} + (1-\gamma)I$,
amplifying hardware noise by factor $\gamma=1,2,3$.
The standard observable, the autocorrelation
$C(t)=\langle Z_c(t)Z_c(0)\rangle$, compresses the entire space-time
evolution into one scalar per cycle, losing spatial information and
dominated by the lightcone interior where $\chi$ is large and the
approximation is least faithful.

Related approaches exploit the lightcone causal structure for error
mitigation: lightcone shading~\cite{eddins2024lightconeshading} tightens
probabilistic error cancellation bounds by restricting to the observable's
causal cone; non-zero noise extrapolation~\cite{thompson2025nonzero}
adds noise to ease tensor network simulation then extrapolates back,
the reverse of our amplification strategy; and Pauli
propagation~\cite{eddins2026noisecanceling} cancels noise classically
before hardware measurement.
Recent work~\cite{chang2026noise} proposes exploiting the coupling
between depolarizing noise and operator weight to extract a global
weight spectrum from noisy circuits at variable noise strength.
We extend this to a site-resolved observable: rather than integrating
over all sites, we track $n(v,t)$ per site per cycle, recovering the
full spatial structure of how noise deforms the lightcone.
The non-identity Pauli weight $n(v,t) = \sum_{P:\,P_v\neq I} c_P(t)^2$
is extracted from the same simulation at no additional cost.
By normalization $\sum_P c_P^2=1$, it is a probability~\cite{nahum2018operator}
and retains the full spatial field at every cycle.
At the lightcone frontier where $\chi\approx1$ and the simulation is
exact, $n(v,t)$ is the most reliable observable the simulation
produces, provided the interior has sufficient $\chi$ to propagate
the signal forward faithfully.
Eigenvalue analysis of the composed PTM confirms both noise models are
in the pure amplitude damping regime (off-diagonal fraction $\approx36\%$,
constant across $\gamma$, $\Delta<0.002$): noise delays and damps the
Hamiltonian-determined wave without redirecting it, making the wavemap
a model-free noise diagnostic.
This is consistent with the analytical result of~\cite{yu2026lightconescalingincircuitnoise}
that a noise event contributes only when it overlaps the Heisenberg
evolution of the measured Pauli operator; our wavemap makes this
overlap structure spatially explicit and site-resolved.
Fitting time-adaptive MPF coefficients $\alpha(t)$ over the frontier
recovers up to 55\% of information loss on heavy-hex and 65\% on
rect $7\times7$; since $\gamma$ maps to the ZNE gate-folding
factor~\cite{giurgica2020digital}, the method applies directly to
real quantum processors.

% ─────────────────────────────────────────────────────────────────────────────
\section{The Pauli Lightcone Wavemap}
\label{sec:wavemap}
This section introduces the wavemap observables following the progression:
\begin{equation}
  C(t) \;\xrightarrow{\text{spatialize}}\; n(v,t)
       \;\xrightarrow{\text{amplify}}\; n_\gamma(v,t)
       \;\xrightarrow{\text{minimize}}\; \min_\alpha \mathcal{L}
  \label{eq:progression}
\end{equation}
\subsection{Autocorrelation as Observable}

The standard observable in tensor network simulation of quantum lattice
models is the autocorrelation $C(t) = \langle Z_c(t) Z_c(0)\rangle$,
computed as an inner product between the evolved operator and the
initial state at the center site $c$.
This compresses the entire space-time evolution into one number per cycle.
Sites strongly affected by noise and sites barely affected contribute
equally; furthermore, the computation is dominated by the interior of
the lightcone, where bond dimension $\chi$ is large and the tensor
network truncation is least faithful.

\subsection{The Pauli Weight $n(v,t)$ as a Spatial Observable}

Any operator $O_c(t)$ initialized at the center and expanded in the
$n$-qubit Pauli basis $\{I,X,Y,Z\}^{\otimes n}$ as
$O_c(t) = \sum_P c_P(t)\,P$ satisfies $\sum_P c_P(t)^2 = 1$.
The non-identity Pauli weight at site $v$ and cycle $t$ is:
\begin{equation}
    n(v,t) = \sum_{P:\,P_v \neq I} c_P(t)^2
    \label{eq:nvt}
\end{equation}
The same simulation that computes $C(t)$ also computes $n(v,t)$ at no
additional cost, retaining the full spatial distribution of operator
support across the lattice at every cycle.


\subsection{The $n(v,t)$ as a Probability}

By normalization, $n(v,t) \in [0,1]$ is the probability that a Pauli
string sampled from the operator distribution acts non-trivially at
site $v$~\cite{nahum2018operator,Tan2025Operator}.
This is not an approximation: it follows directly from $\sum_P c_P^2=1$.
The field $n(v,t)$ is therefore a probability distribution over the
lattice, evolving in time as the operator spreads outward from the
initialization site.

\subsection{Lightcone: Velocity and Information}

The \emph{frontier} at cycle $t$ is the set of sites first crossing a
significance threshold $\varepsilon_{\rm th} = 10^{-4}$:
\begin{equation}
    \mathcal{F}(t) = \{v : n(v,t) > \varepsilon_{\rm th},\;
                          n(v,t-1) \leq \varepsilon_{\rm th}\}
\end{equation}
Its maximum breadth-first search (BFS) distance grows linearly at
velocity $v_{\rm hop}$ (hops/cycle), establishing the Lieb-Robinson
cone.
The wavemap observables are computed at the leading edge of the
lightcone where they are least sensitive to bond dimension truncation.
The threshold $\varepsilon_{\rm th}$ plays the role of a physical
tolerance — analogous to the truncation cutoff $\delta$ in the
simulation — defining the precision with which we observe the
wavefront's presence and amplitude ratio.
The choice of $\chi$ reflects three practical constraints: memory,
wall-clock time, and accuracy; its effect on $C(t)$ convergence
and collapse was studied separately in~\cite{dalberto2026ising}.
A full characterization of the interplay between $\chi$,
$\varepsilon_{\rm th}$, and wavemap accuracy is left for future work.
Crucially, the simulation captures non-uniform frontiers naturally:
irregular topology, per-bond coupling heterogeneity, and asymmetric
noise all shape $\mathcal{F}(t)$ without any model assumption, as
the heavy-hex results demonstrate (near-degenerate $v_{\gamma_2}
\approx v_{\gamma_3}$, non-monotone delay, Table~\ref{tab:vhop}).

\subsection{Noiseless and Noisy Waves}

We simulate four variants: noiseless (NL) and three noise amplification
levels $\gamma=1,2,3$ (chosen from the ZNE literature~\cite{giurgica2020digital}),
where $\gamma$ rescales the noise channel $M_\gamma = \gamma M_{\rm hw} + (1-\gamma)I$.
Each variant produces its own wavefront advancing at velocity
$v_{\gamma} \leq v_{\rm nl}$: noise suppresses amplitudes below
threshold, delaying the frontier crossing.
For each site $v$ and noise level $\gamma$ we define the
\emph{arrival delay} $l_\gamma(v) = t_\gamma(v) - t_{\rm nl}(v) \geq 0$
and the \emph{amplitude ratio}
$r_\gamma(v) = n_\gamma(v,t_\gamma(v))/n_{\rm nl}(v,t_\gamma(v)) \leq 1$,
both guaranteed by the Lieb-Robinson bound for pure amplitude damping.
Together they form the \emph{wavemap}: a spatial portrait of how noise
deforms the lightcone, site by site and cycle by cycle.

\subsection{Cross-Entropy as the Information Cost}

Since $n(v,t)$ is a probability, the natural measure of deviation
between the noisy and noiseless is by loss of information.  As
Jaynes~\cite{Jaynes1957PartI,Jaynes1957PartII} established, when
quantities are probabilities, information-theoretic objectives are the
correct and unifying choice: MSE penalizes amplitude error uniformly,
while cross-entropy penalizes information loss proportionally to its
magnitude.  We define the per-site cross-entropy loss summed over a
delay-adaptive window $[t_\gamma(v)-\Delta,\,t_\gamma(v)+\Delta]$,
$\Delta = \max(1,l_\gamma(v))$:
\begin{equation}
    \mathcal{L}_\gamma(v) = \sum_{t \in \text{window}}
    -\ln \frac{n_\gamma(v,t)}{n_{\rm nl}(v,t)}
    \label{eq:loss}
\end{equation}
Zero when the distributions coincide, positive when noise suppresses
probability, with window width set by the delay itself so no free
parameter is introduced.  This is the accumulated information loss at
site $v$: the wavemap observable that quantifies, site by site, how
much the noisy evolution has forgotten relative to the noiseless
reference.

The next section shows how these observables enable MPF error recovery.

% ─────────────────────────────────────────────────────────────────────────────
\begin{table*}[!t]
\centering
\setlength{\tabcolsep}{3pt}
\small
\begin{tabular}{r|rrrr|rrrr}
\toprule
 & \multicolumn{4}{c|}{Heavy-hex $3\times3$} & \multicolumn{4}{c}{Rect $7\times7$} \\
$t$ & NL & $\gamma_1$ & $\gamma_2$ & $\gamma_3$ & NL & $\gamma_1$ & $\gamma_2$ & $\gamma_3$ \\
\midrule
 1 & $-$.924 & $-$.924 & $-$.925 & $-$.925 & $-$.914 & $-$.916 & $-$.918 & $-$.920 \\
 2 & $+$.832 & $+$.836 & $+$.839 & $+$.843 & $+$.827 & $+$.843 & $+$.857 & $+$.870 \\
 3 & $-$.763 & $-$.773 & $-$.783 & $-$.792 & $-$.748 & $-$.791 & $-$.826 & $-$.854 \\
 4 & $+$.789 & $+$.803 & $+$.816 & $+$.828 & $+$.612 & $+$.696 & $+$.759 & $+$.808 \\
 5 & $-$.828 & $-$.848 & $-$.864 & $-$.878 & $-$.500 & $-$.623 & $-$.711 & $-$.774 \\
 6 & $+$.851 & $+$.876 & $+$.895 & $+$.909 & $+$.414 & $+$.574 & $+$.685 & $+$.761 \\
 7 & $-$.838 & $-$.870 & $-$.892 & $-$.907 & $-$.341 & $-$.534 & $-$.665 & $-$.751 \\
 8 & $+$.817 & $+$.857 & $+$.882 & $+$.897 & $+$.287 & $+$.508 & $+$.654 & $+$.744 \\
 9 & $-$.804 & $-$.853 & $-$.881 & $-$.896 & $-$.248 & $-$.495 & $-$.650 & $-$.742 \\
10 & $+$.793 & $+$.853 & $+$.883 & $+$.898 & $+$.219 & $+$.490 & $+$.649 & $+$.740 \\
\bottomrule
\end{tabular}
\caption{$C(t)$ for NL and $\gamma=1,2,3$, 10 cycles ($\varepsilon_{\rm th}=10^{-4}$).
Heavy-hex: noise levels nearly indistinguishable.
Rect $7\times7$: monotone decay, spatial information lost.
$\chi$ grows $4\to163$ on heavy-hex; saturates at $\chi=50$ on rect $7\times7$
(truncation error present at interior sites; wavemap observables computed at
frontier where $\chi\approx1$ and simulation is exact.)}
\label{tab:ct}
\end{table*}

\section{Results}
\label{sec:results}

Following the progression of Eq.~\ref{eq:progression}, we present
results in three stages: the autocorrelation $C(t)$, the frontier
velocity $v_{\rm hop}$ from $n(v,t)$, and the MPF entropy recovery.

\subsection{Stage 1: Autocorrelation $C(t)$}

Existing MPF methods operate at two levels of the same scalar observable.
The first fits $C(t)$ directly: $\sum_i c_i C_{\gamma_i}(t) \approx C_{\rm nl}(t)$,
one equation per cycle.
The second solves the bilinear system $A\alpha = B$~\cite{mpf-ref,robertson2026lindbladmultiproductformulas}, where
$A_{ij} = \langle\psi_{\gamma_i}|\psi_{\gamma_j}\rangle$ are noisy-noisy
inner products and $B_i = \langle\psi_{\gamma_i}|\psi_{\rm nl}\rangle$
are noisy-noiseless inner products: a richer fit, but still a global
scalar per cycle, still dominated by the interior of the lightcone
where $\chi$ is large and the approximation is least faithful.

Table~\ref{tab:ct} shows $C(t)$ at each cycle for NL and
$\gamma=1,2,3$ on both topologies.
On heavy-hex, $C(t)$ oscillates and the noise levels are nearly
indistinguishable: the scalar observable is almost blind to noise.
On rect $7\times7$, $C(t)$ decays monotonically and the $\gamma$
levels separate, but all spatial structure is lost.
In both cases, $C(t)$ and $A\alpha=B$ compress the entire
space-time evolution into one number per cycle.

\subsection{Stage 2: Frontier Velocity from $n(v,t)$}

The non-identity Pauli weight $n(v,t)$ resolves the spatial structure
invisible to $C(t)$.
Table~\ref{tab:vhop} shows the frontier velocity $v_{\rm hop}$
(hops/cycle), the slope of the maximum BFS frontier distance versus cycle.

\begin{table}[t]
\centering
\setlength{\tabcolsep}{2pt}
\scriptsize
\begin{tabular}{llr rrr rrr rrr}
\toprule
 & & & \multicolumn{3}{c}{$\gamma=1$} & \multicolumn{3}{c}{$\gamma=2$} & \multicolumn{3}{c}{$\gamma=3$} \\
Topology & $v_{\rm nl}$ & $n_{\rm nl}$ & $n$ & $\mu_l$ & $\sigma_l$ & $n$ & $\mu_l$ & $\sigma_l$ & $n$ & $\mu_l$ & $\sigma_l$ \\
\midrule
HH $3\times3$ & 0.376 & 21 & 8 & 1.13 & 1.97 & 10 & 2.10 & 2.81 & 8 & 0.75 & 1.30 \\
Rect $7\times7$ & 0.448 & 49 & 49 & 0.98 & 1.15 & 37 & 1.35 & 1.36 & 25 & 1.24 & 1.77 \\
\bottomrule
\end{tabular}
\caption{Frontier velocity $v_{\rm nl}$ and delay statistics
$(n,\mu_l,\sigma_l)$ per $\gamma$: sites reached, mean and std of
arrival delay $l_\gamma(v)$ in cycles.}
\label{tab:vhop}
\end{table}

The frontier velocity $v_{\rm hop}$ is not a property of the noise alone.
It is a joint function of three factors: the graph topology (minimum hop
distance between connected sites, coordination number, girth), the
Hamiltonian coupling $J_{ij}$ (which determines how much Pauli weight
transfers per gate application), and the noise representation at gate
level (how $\gamma$ rescales the effective amplitude per bond crossing).
The per-site delay statistics $(n_{\rm reached},\,\mu_l,\,\sigma_l)$
in Table~\ref{tab:vhop} provide the statistical fingerprint of all
three combined: sites that the noisy wave fails to reach entirely
($n_{\rm reached} < n_{\rm nl}$), the mean delay for those it does
reach ($\mu_l$), and the spread ($\sigma_l$) reflecting path-dependent
variation in coupling and noise accumulation.
The wavemap therefore characterizes the interplay of geometry, dynamics,
and noise in a single measurement.

Sites outside the noiseless lightcone have $\chi=1$ identically: the
operator has not yet reached them and they carry no correlations.
Simulation outside the cone is unnecessary: the arrival time
$t_{\rm nl}(v) = d(c,v)/v_{\rm hop}$ provides a per-site entry
schedule, enabling lightcone-aware memory allocation that eliminates
computation before each site activates.
Crucially, the wavemap turns this boundary into a diagnostic: any
$n_\gamma(v,t) > \varepsilon_{\rm th}$ observed \emph{outside} the
noiseless cone is not silence but a signal: spontaneous activation
that the frontier flags as anticausal.
For pure amplitude damping this never occurs ($l_\gamma(v) \geq 0$
always); its presence is a fingerprint of rotation-type noise injecting
Pauli weight independently of the Hamiltonian causal structure.

\subsection{Stage 3: MPF Entropy Recovery}

The wavemap cross-entropy loss $\mathcal{L}_\gamma(v)$ is both a
diagnostic and an objective: minimizing it over MPF coefficients
recovers the noiseless Pauli weight field from noisy samples.

All methods reconstruct $n_{\rm MPF}(v,t) = \sum_i c_i\, n_{\gamma_i}(v,t)$
and are evaluated by the global mean cross-entropy loss $\mu_\mathcal{L}$
(Eq.~\ref{eq:loss}), lower is better.

A geometric baseline fits $c$ by minimizing squared error on the
frontier distance series $\min_c \sum_t (F_{\rm MPF}(t)-F_{\rm nl}(t))^2$,
subject to $F_{\rm MPF}(t)\leq F_{\rm nl}(t)$.
This objective sees only scalar distances, never amplitudes.
Evaluating its $n_{\rm MPF}$ on $\mu_\mathcal{L}$ yields 1.058/3.911
with no improvement over the best noisy sample, confirming that
geometric and information-theoretic objectives are not aligned.

\textbf{MPF entropy} fits $c$ by directly minimizing $\mu_\mathcal{L}$:
\begin{equation}
    \min_{c_i}\; \sum_{v,t} -\ln\frac{n_{\rm MPF}}{n_{\rm nl}},\quad
    \sum_i c_i = 1,\quad n_{\rm MPF} \leq n_{\rm nl}
    \label{eq:entropy_mpf}
\end{equation}
Table~\ref{tab:entropy} shows $\mu_\mathcal{L}$ for the noisy samples
and MPF entropy on both topologies at 10 cycles; the geometric baseline
is excluded as it optimizes a different objective.

\begin{table}[t]
\centering
\small
\begin{tabular}{lcc}
\toprule
Method & $\mu$ hh dep & $\mu$ rect $7{\times}7$ \\
\midrule
$\gamma=1$       & 1.058 & 3.680 \\
$\gamma=2$       & 1.538 & 7.432 \\
$\gamma=3$       & 2.277 & 8.003 \\
\midrule
MPF entropy      & 0.596 & 2.783 \\
\bottomrule
\end{tabular}
\caption{Mean per-site $\mu_\mathcal{L}$ (lower is better), 10 cycles.
MPF entropy reduces loss 44\% on heavy-hex and 24\% on rect $7\times7$.}
\label{tab:entropy}
\end{table}

\subsection{Time-Adaptive Coefficients $\alpha(t)$}

The global MPF entropy fit uses fixed $c$ across all space and time.
We refine this by fitting one coefficient vector $\alpha(t)\in\mathbb{R}^3$
per cycle over the union frontier
$\mathcal{F}(t) = \mathcal{F}_{\rm nl}(t) \cup \bigcup_i \mathcal{F}_{\gamma_i}(t)$,
a thin spatial shell where $\chi\approx1$ in all simulations and every
sample is exact.
The approximation is in the application: coefficients fitted on the
frontier are extrapolated to all sites at that cycle.

The dual, $\alpha(d)$, fixes a BFS shell and fits over a time window,
a $\Delta$ in time rather than space.
Table~\ref{tab:lt} shows $L(t) = \sum_v -\ln(n_{\rm MPF}/n_{\rm nl})$,
the global cross-entropy loss summed over active sites at cycle $t$, for $\gamma=1$, $\alpha(t)$,
and $\alpha(d)$: heavy-hex at 20 cycles (to expose the frontier
collapse at cycles 12--13), rect $7\times7$ at 10 cycles.
Table~\ref{tab:ld} shows $L(d)$ per BFS shell at 10 cycles for both
topologies.

\begin{table}[t]
\centering
\setlength{\tabcolsep}{3pt}
\scriptsize
\begin{tabular}{r|rrr|rrr}
\toprule
 & \multicolumn{3}{c|}{Heavy-hex} & \multicolumn{3}{c}{Rect $7\times7$} \\
$t$ & $\gamma_1$ & $\alpha(t)$ & $\alpha(d)$ & $\gamma_1$ & $\alpha(t)$ & $\alpha(d)$ \\
\midrule
 1 &  0.11 &  0.00 &  0.00 &  0.85 &  0.01 &  0.01 \\
 2 &  0.23 &  0.02 &  0.00 &  3.17 &  0.25 &  0.66 \\
 3 &  0.91 &  0.16 &  0.00 &  8.83 &  1.56 &  2.00 \\
 4 &  1.33 &  0.29 &  0.00 & 20.00 &  5.42 & 18.88 \\
 5 &  1.61 &  0.42 &  0.00 & 33.51 & 12.16 & 39.20 \\
 6 &  3.58 &  1.22 &  2.16 & 43.47 & 18.45 & 53.56 \\
 7 &  5.30 &  2.18 &  4.90 & 44.96 & 19.46 & 55.85 \\
 8 &  6.76 &  3.43 &  6.86 & 46.06 & 20.64 & 57.51 \\
 9 &  8.11 &  4.89 &  8.03 & 47.41 & 21.69 & 59.26 \\
10 &  8.58 &  5.44 &  8.47 & 49.06 & 23.06 & 61.16 \\
11 & 14.90 & 14.08 & 18.33 & -- & -- & -- \\
12 & 22.12 & 26.26 & 29.75 & -- & -- & -- \\
13 & 23.51 & \textit{41.83} & 31.81 & -- & -- & -- \\
14 & 24.69 & 28.76 & 33.47 & -- & -- & -- \\
15 & 27.60 & 32.36 & 37.88 & -- & -- & -- \\
16 & 28.65 & \textbf{7.35} & 39.38 & -- & -- & -- \\
17 & 38.48 & 12.77 & 53.10 & -- & -- & -- \\
18 & 39.36 & 14.34 & 54.15 & -- & -- & -- \\
19 & 42.89 & 16.72 & 58.87 & -- & -- & -- \\
20 & 50.29 & 21.62 & 68.60 & -- & -- & -- \\
\midrule
$\mu$ & 1.058 & \textbf{0.481} & 1.050 & 3.680 & \textbf{1.276} & 5.808 \\
\bottomrule
\end{tabular}
\caption{$L(t) = \sum_v -\ln(n_{\rm MPF}/n_{\rm nl})$: global cross-entropy loss per cycle.
Hh at 20 cycles, r7 at 10 (-- = not run).
\textit{Italics}: $\alpha(t)$ ill-conditioned (frontier collapse, cycles 12--15).
\textbf{Bold}: recovery after frontier repopulates.
$\mu$ row: mean $\mu_\mathcal{L}$. Values floored at zero.\footnote{Not all
per-site Lieb-Robinson constraints enforced; a full pointwise solver would
eliminate floored entries.}}
\label{tab:lt}
\end{table}

\begin{table}[htb]
\centering
\setlength{\tabcolsep}{3pt}
\scriptsize
\begin{tabular}{r|rrr|rrr}
\toprule
 & \multicolumn{3}{c|}{Heavy-hex} & \multicolumn{3}{c}{Rect $7\times7$} \\
$d$ & $\gamma_1$ & $\alpha(t)$ & $\alpha(d)$ & $\gamma_1$ & $\alpha(t)$ & $\alpha(d)$ \\
\midrule
0 & 0.00 & 0.00 & 0.00 & 0.00 & 0.00 & 0.00 \\
1 &    0.03 & 0.00 &    0.00 &    0.07 & 0.00 &    0.00 \\
2 &    0.07 & 0.01 &    0.00 &    0.25 & 0.02 &    0.01 \\
3 &    0.21 & 0.05 &    0.00 &    0.62 & 0.13 &    0.10 \\
4 &    0.51 & 0.24 &    0.92 &    1.11 & 0.40 &    1.79 \\
5 &    0.88 & 0.62 &    1.50 &    1.70 & 0.84 &    2.53 \\
6 &     --  &  --  &    --   &    2.33 & 1.38 &    3.27 \\
\bottomrule
\end{tabular}
\caption{Mean cross-entropy loss $\mathcal{L}(d)$ per BFS shell.
$\alpha(t)$ reduces loss at every shell; $\alpha(d)$ degrades at
outer shells (too few sites per shell to constrain the fit).}
\label{tab:ld}
\end{table}

\subsection{Connection to Hardware: ZNE}

The noise amplification parameter $\gamma$ maps directly to the
gate-folding factor in zero-noise extrapolation
(ZNE)~\cite{giurgica2020digital}: running the circuit at fold factors
$\gamma=1,2,3$ multiplies the effective noise per gate by $\gamma$,
exactly the $M_\gamma$ model above.
The Pauli weight $n(v,t)$ is accessible per qubit per cycle via local
Pauli measurements.
On real quantum hardware there is no bond dimension: the full quantum
state is available exactly, and $n(v,t)$ is measured directly without
truncation artifact. The frontier is then the leading edge of the
measured operator support, defined by the same threshold
$\varepsilon_{\rm th}$ but exact at every site.
The wavemap observables and $\alpha(t)$ fitting then follow from the
same computation applied to experimental data, no simulation required.
The present work provides the ground truth and the methodology;
hardware validation is a direct next step.

% ─────────────────────────────────────────────────────────────────────────────
\section{Conclusion}
\label{sec:conclusion}
A brief list of consideration follows: First, eigenvalue analysis of
the composed gate-plus-noise PTM confirms pure amplitude damping: the
spatial propagation pattern is entirely a property of the Hamiltonian,
validating the wavemap as a model-free noise diagnostic.  Second, the
Lieb-Robinson constraint $n_{\rm MPF}\leq n_{\rm nl}$ is a hard causal
bound, not a regularizer; enforcing it pointwise is an open problem
whose solution would eliminate the floored entries in
Table~\ref{tab:lt}.  Third, fitting time-adaptive coefficients
$\alpha(t)$ over the union frontier recovers 55\% of information loss
on heavy-hex and 65\% on rect $7\times7$, outperforming both fixed
global coefficients and space-adaptive $\alpha(d)$. Fourth, fast
GPU-accelerated simulation enables systematic exploration of the
causal relations between noise character, Hamiltonian coupling, and
wavefront structure across a library of noise models and topologies.
This opens the path to classifiers trained on simulation ground truth
that diagnose noise type and severity directly from measured $n(v,t)$
on real quantum processors, closing the loop between classical
simulation and hardware characterization.

\bibliography{qiml26}

\end{document}
